% Einheit 3 von 4 – Der Graphenblick (bank/stammfunktion-und-hauptsatz/e3.jsonl, zusammenbau v0.1)
%% TODO Zeitmarke Einheit 3: Marken-Zeile nicht lesbar
%% TODO Zweigzeile Teil 1: was hier gelernt wird (Satz in Schülersprache aus dem Einheitstitel) – Einheit 3
\einheitenkopf[e3]{Einheit 3 von 4 · Der Graphenblick}
\zweigzeile{Abitur GK}
\setcounter{aufgabe}{12}

%% TODO Titel: Ich-kann-Satz fehlt in der Bank (Platzhalter: Kettenname) – Nr. 13
%% TODO Anweisung: ein Satz über der ersten Teilaufgabe fehlt in der Bank – Nr. 13
\begin{aufgabe}{Der Graphenblick}
%% TODO \rechenplatz{n}: Schrittzahl des Grundfalls fehlt in der Bank (nur bei fester Schreibform, 2.3 b)
\begin{teile}
\teil $u(x) = x^3 + 2$ und $v(x) = 3x^2$. Eine der beiden ist Stammfunktion der anderen. Welche ist die Stammfunktion? \\ \kreuz{$u$} \\ \kreuz{$v$}
\teil Das Bild zeigt den Graphen von $f$. Skizziere den Graphen der Stammfunktion $F$ von $f$, die durch $P(0 | 1)$ geht.

\begin{ksys}[xmin=-2,xmax=4,ymin=-3,ymax=4] \gerade{1}{-1}{f} \punkt{0}{1}{P} \end{ksys}
\teil Das Bild zeigt den Graphen einer Stammfunktion $F$ von $f$. Lies passende Werte ab und berechne das Integral von $2$ bis $6$ über $f(x)$. \hfill (Abitur 2022 LK) \\

\begin{ksys}[xmin=0,xmax=8,ymin=-2,ymax=6,ablesen] \funktionab{2-(\x-4)^3/8}{F}{1}{7} \end{ksys}
Wert: \leerfeld
\teil Das Bild zeigt den Graphen einer Stammfunktion $F$ von $f$. Zeichne die Tangente an den Graphen von $F$ bei $x = 2$ ein und bestimme damit $f(2)$.

\begin{ksys}[xmin=-2,xmax=4,ymin=-2,ymax=4] \parabel{0.5}{1}{-0.5}{F} \end{ksys}
$f(2) = $ \leerfeld
\teil Begründe, dass jede Stammfunktion von $g(x) = x^2 \cdot (x - 2)$ bei $x = 2$ ihr Minimum annimmt. \hfill (Abitur 2026 GK)
\teil Das Bild zeigt den Graphen von $g$ mit den Nullstellen $-3$, $1$ und $4$. Aussage: Der Graph jeder Stammfunktion von $g$ hat zwei Tiefpunkte. Beurteile die Aussage. \hfill (Abitur 2022 LK) \\ \kreuz{wahr} \\ \kreuz{falsch}

\begin{ksys}[xmin=-4,xmax=5,ymin=-5,ymax=6,ablesen] \funktionab{0.25*(\x+3)*(\x-1)*(\x-4)}{g}{-3.5}{4.7} \end{ksys}
\teil Das Bild zeigt den Graphen von $f'$. Gib zwei verschiedene Zahlen $a$ und $b$ an, für die das Integral von $a$ bis $b$ über $f''(x)$ null ist.

\begin{ksys}[xmin=-3,xmax=5,ymin=-3,ymax=4,ablesen] \parabel{0.5}{1}{-2}{f'} \end{ksys}
$a = $ \leerfeld , $b = $ \leerfeld
\end{teile}
\end{aufgabe}

%% TODO Titel: Ich-kann-Satz fehlt in der Bank (Platzhalter: Kettenname) – Nr. 14
%% TODO Anweisung: ein Satz über der ersten Teilaufgabe fehlt in der Bank – Nr. 14
\begin{aufgabe}{Der Graphenblick – Fehler finden, Begründen}
\begin{teile}
\teil Ich finde den Fehler: $f(x) = x^2 - 4$ hat bei $x = 0$ einen Tiefpunkt. Jonas schreibt: Also hat jede Stammfunktion von $f$ bei $x = 0$ einen Tiefpunkt.
\teil Begründe, warum eine Stammfunktion $F$ von $f$ an einer Nullstelle von $f$ ohne Vorzeichenwechsel keinen Extrempunkt hat.
\end{teile}
\end{aufgabe}
